\vspace{-1mm}
\section{Related Work}
\label{sec:relatedwork}
\vspace{-1mm}
\subsection{Image Animation}
\label{subsec:image_animation}
Generating animation from still images is a heavily studied research area. Early physical simulation-based approaches~\cite{jhou2015animating,chuang2005animating} focus on simulating the motion of specific objects, % with known or computed identities, 
resulting in low generalizability due to the independent modeling of each object category. To produce more realistic motion, reference-based methods~\cite{prashnani2017phase,okabe2009animating,siarohin2019first,cheng2020time,karras2023dreampose,siarohin2019animating,siarohin2021motion,wang2021latent} transfer motion or appearance information from reference signals, such as videos, to the synthesis process. Although they demonstrate better temporal coherence, the need for additional guidances limits their practical application. Additionally, a stream of works based on GAN~\cite{shaham2019singan,hinz2021improved,karras2020analyzing} can generate frames by perturbing initial latents or performing random walk in the latent vector space. However, the generated motion is not plausible since the animated frames are just a visualization of the possible appearance space without temporal awareness. Recently, (learned) motion prior-based methods~\cite{endo2019animating,holynski2021animating,mallya2022implicit,ni2023conditional,weng2019photo,zhao2022thin,xiao2023automatic,hu2022make} animate still images through explicit or implicit image-based rendering with estimated motion field or geometry priors. Similarly, video prediction~\cite{zhang2020dtvnet,xiong2018learning,babaeizadeh2018stochastic,li2018flow,xue2016visual,voleti2022mcvd,hu2023dynamic,hoppe2022diffusion,franceschi2020stochastic} predicts future video frames starting from single images by learning spatio-temporal priors from video data.


Although existing approaches has achieved impressive performance, they primarily focus on animating motions in curated domains, particularly stochastic~\cite{endo2019animating,jhou2015animating,blattmann2021ipoke,xue2018visual,okabe2009animating,chuang2005animating,lee2018stochastic,dorkenwald2021stochastic} and oscillating~\cite{li2023generative} motion. Furthermore, the animated objects are limited to specific categories, \eg, fluid~\cite{holynski2021animating,okabe2009animating,holynski2021animating,mahapatra2022controllable}, natural scenes~\cite{xiong2018learning,li2023generative,cheng2020time,jhou2015animating,shaham2019singan}, human hair~\cite{xiao2023automatic}, portraits~\cite{wang2021latent,wang2020imaginator,geng2018warp}, and bodies~\cite{wang2021latent,weng2019photo,karras2023dreampose,siarohin2021motion,blattmann2021understanding,bertiche2023blowing}. In contrast, our work proposes a generic framework for animating open-domain images with a wide range of content and styles, which is extremely challenging due to the overwhelming complexity and vast diversity.


\vspace{-0.5mm}
\subsection{Video Diffusion Models}
\label{subsec:video_diffusion_models}
\vspace{-0.5mm}
Diffusion models (DMs)~\cite{sohl2015deep,ho2020denoising} have recently shown unprecedented generative power in text-to-image (T2I) generation~\cite{ramesh2022hierarchical,nichol2022glide,saharia2022photorealistic,rombach2022high,he2023scalecrafter,zhang2023real}. To replicate this success to video generation, the first video diffusion model (VDM)~\cite{ho2022video} is proposed to model low-resolution videos using a space-time factorized U-Net in pixel space. Imagen-Video~\cite{ho2022imagenvideo} presents effective cascaded DMs with $\mathbf{v}$-prediction for generating high-definition videos. To reduce training costs, subsequent studies~\cite{wang2023modelscope,he2022latent,blattmann2023align,zhou2022magicvideo,wang2023lavie} are engaged in transferring T2I to text-to-video (T2V)~\cite{singer2022make,luo2023videofusion,ge2023preserve,zhang2023show1}, and learning VDMs in latent or hybrid-pixel-latent space.

Although these models can generate high-quality videos, they only accept text prompts as the sole semantic guidance, which can be vague and may not accurately reflect users' intention. Similar to adding controls in T2I~\cite{zhang2023adding,mou2023t2i,ye2023ipadapter,shi2023instancebooth}, introducing control signals in T2V, such as structure~\cite{esser2023structure,xing2023make}, pose~\cite{ma2023follow,zhang2023controlvideo}, and Canny edge~\cite{khachatryan2023text2video}, has been increasingly receiving much attention.
However, visual conditions in VDMs~\cite{tang2023any,yin2023dragnuwa}, such as RGB images, remain under-explored. Most recently and concurrently, image condition is examined in Seer~\cite{gu2023seer}, VideoComposer~\cite{wang2023videocomposer},  and I2VGen-XL~\cite{I2VGen-XL} for (text-)image-to-video synthesis.
However, they either focus on the curated domain, \ie, indoor objects~\cite{gu2023seer}, or fail to generate temporally coherent frames and realistic motions~\cite{wang2023videocomposer} and preserve visual details of the input image~\cite{I2VGen-XL} due to insufficient context understanding and loss of information of the input image. 
Moreover, recent proprietary T2V models~\cite{singer2022make,molad2023dreamix,villegas2023phenaki,yu2023magvit} have been demonstrated to be extensible to image-to-video synthesis. However, their results rarely adhere to the input image and suffers from the unrealistic temporal variation issue. Our approach is built upon text-conditioned VDMs to leverage their rich dynamic prior for animating open-domain images, by incorporating tailored designs for better semantic understanding and conformity to the input image.